\scanpage{9}
\deleted{\uline{Corollaire}}

\uline{Proposition 3} Soit $A$ un schéma abélien sur $V$,
et $u':A':A_K\longrightarrow B'$ une isogénie, où $B'$ est
un schéma abélien sur $K$. Il existe alors
un schéma abélien $B$ sur $V$, \deleted{\unsure{se réduisant à} \unsure{un}}
tel que $B\otimes_V K\simeq B'$, et un \deleted{\unsure{en plus|et cette} isogénie}
ho\inserted{\unsure{mo}}morphisme \deleted{$A\mathrel{-}$} de schémas abéliens $u:A\mathrel{-}B$
induisant $u':A'\mathrel{-}B'$.

N.B. On ne se sert que du th. de Weil,
disant qu'un $V$-morphisme \deleted{\unsure{$A\mathrel{-}B$|$A_K\to B_K$} \unsure{entre}}
$A'\mathrel{-}B'$ \unsure{provient} d'un morphisme $A\mathrel{-}B$.

\uline{Dém}. Soit $H'$ le noyau de $u'$,
c'est un sous-groupe fini de $A'$, donc
provient d'un sous-groupe $V$-plat $H=\overline{H'}$ de
$A$, nécessairement \uline{fini} sur $V$, \unsure{étant} propre \inserted{et plat} sur
% slanted
\margin{\unsure{référence|retenons}.}$V$ à fibre générale finie. Comme
$A$ est \uline{projectif} sur $V$ (\unsure{réf}), on
peut construire le \inserted{schéma} quotient $A/\unsure{\Gamma|H}$,
qui est un schéma en groupes
\unsure{car|sur} $A$ est plat sur $V$, \uline{plat} sur $V$,
($=$ raisons), propre sur $V$, et tel que
% archive sheet number: 8
